% Integrity verification methods using Message Authentication Codes (MAC) often impose significant computational and bandwidth overhead. While the short MAC was initially proposed to reduce this overhead, it compromises security. Alternative schemes, such as aggregate, compound, and progressive MACs, have been developed to maintain security levels while reducing overhead and computational demands. However, these heuristic methods remain susceptible to Denial of Service (DoS) and jamming attacks, particularly in critical wireless environments. This paper develops a new optimization framework for integrity mechanism that is based on a novel mapping of the problem to a graph, allowing for the deployment of mathematical optimizations. OptiMAC maximizes the expected number of tag bits per message, ensuring optimal security measures for integrity check system under adversarial scenarios. This is achieved by our novel unifying framework and innovative metrics to represent security and efficiency, enabling the optimization model to identify optimal message-tag assignments. Through extensive real-world testing on live webcam video streaming over Wi-Fi and 5G, we demonstrate a constant improvement in the number of tag bits per message compared to existing methods, validating our methodology. 
% This paper presents an optimized MAC aggregation framework that enhances integrity checks in wireless communications under adversarial conditions. Building on the Progressive MAC (ProMAC) model, our approach systematically defines parameters like tag generation and update rules, allowing adaptive optimization for resilience against forgery, denial of service (DoS), and lossy channel conditions.
% Leveraging existing security proofs from ProMAC and aggregate MAC studies, our framework introduces the Strength Number (SN) metric to enhance resilience. By optimizing the tag-to-message ratio within a dependency area calculated as \(2^{n \times m}\), we create a broad solution space that improves both latency and security.
% Experimental results across varying attack rates show that higher dependency areas (e.g., a 15/30 tag-to-message ratio) outperform configurations with similar ratios (like 10/20) in terms of SN and latency. This scalable framework is particularly suited for secure, resource-constrained wireless applications, demonstrating an adaptable, efficient approach to improving MAC scheme reliability in hostile environments.
Integrity verification utilizing Message Authentication Codes (MACs) can impose high overhead by generating and transmitting extra bits. Alternatively, grouping multiple messages and sending only one MAC for the group can reduce the computational and energy expenditure. However, in such grouping scenarios, the MAC depends on the each of the messages to validate the integrity of the entire group. Therefore, even with a single message loss, the MAC becomes unusable. In adversarial wireless environments, such as under jamming, this grouping strategy can make the system vulnerable to a denial-of-service attack. To address this, we propose OptiMAC, which addresses the shortcomings of the Progressive MAC (ProMAC) framework to increase the resiliency to message losses. In the ProMAC framework, each MAC is derived from an evolving state that holds a buffer of multiple messages and is updated with every new message. While the ProMAC framework's security is proven in principle, it lacks a systematic approach to defining the optimal state update rule for varying channel quality. Our proposed OptiMAC framework (i) formulates the state update process as a message grouping optimization problem and maximizes the number of usable MACs for any jamming probability, and (ii) uses simulations to find the optimum MAC length, which further maximizes throughput. Experiments over WiFi and 5G indicate improvements in security and efficiency.
